\chapter{守恒律：Noether 定理的全面展开}
\label{ch:conservation}

\section{概述}

R2（最小作用量原理）最优雅的推论是 Noether 定理。Emmy Noether 在 1918 年证明：\textbf{作用量的每一个连续对称性都对应一个守恒量。}

"对称性"在日常生活中指"某种变换下不变"。例如，一个完美圆球绕轴旋转，它看起来完全相同——这就是旋转对称。在物理学中，对称性指的是：如果你对整个实验做某种变换（比如移动到一个新位置、旋转一个角度、或者推迟一段时间），物理定律的形式保持不变。

本章展示：
\begin{itemize}
  \item 时间平移不变 $\to$ 能量守恒。这意味着昨天和今天物理定律一样——因此总能量不变
  \item 空间平移不变 $\to$ 动量守恒。北京和纽约的物理定律一样——因此总动量不变
  \item 旋转不变 $\to$ 角动量守恒。朝北和朝东的物理定律一样——因此角动量不变
  \item $U(1)$ 规范不变 $\to$ 电荷守恒。电子的全局相位任意旋转——因此电荷不变
\end{itemize}

这不是巧合。守恒律是宇宙几何对称性的直接表现。如果明天物理定律突然改变（时间平移对称破缺），能量就不再守恒。

\section{Noether 定理的核心思想}

\begin{lawbox}[Noether 定理]
  如果作用量 $S = \int L\,dt$ 在某个连续变换下不变，则存在一个对应的守恒量 $Q$，满足 $dQ/dt = 0$。

  \textbf{前提}：R2（$\delta S = 0$）+ 连续对称性

  \textbf{通俗理解}：想象你在一个完全均匀的空间中——每个点都一样。你把实验设备从 A 点搬到 B 点，实验结果完全相同。Noether 定理说，这种"每个点都一样"的均匀性必然导致某个量守恒。具体是什么量守恒，取决于"什么变换下不变"——时间是动量，旋转是角动量。
\end{lawbox}

\section{能量守恒}

\begin{lawbox}[Energy Conservation]
\[
\frac{dE}{dt} = 0 \quad\text{（孤立系统内总能量不随时间变化）}
\]
更精确的局域形式（每一点的能量流动）：$\partial_\mu T^{\mu\nu} = 0$。

\textbf{前提}：Noether 定理 + 时间平移对称（拉格朗日量 $L$ 不显含时间）

\textbf{变量}：
\begin{itemize}
  \item $E$ — 总能量（J，焦耳）。包括动能（运动能量）、势能（位置能量）、热能（分子随机运动的能量）、化学能等所有形式的能量之和
  \item $T^{\mu\nu}$ — 应力-能量张量。$T^{00}$ 是能量密度（每立方米有多少焦耳），$T^{0i}$ 是能量流动的速率
  \item $dE/dt = 0$ — "变化率为零"意味着这个量是常数。不管系统内部发生什么变化，总能量始终不变
\end{itemize}
\end{lawbox}

\begin{derivationbox}[推导：时间平移不变 → 能量守恒]
  \textbf{第 1 步：什么是"时间平移不变"？}
  "拉格朗日量 $L$ 不显含时间"的数学表达是 $\partial L/\partial t = 0$。物理上，这意味着物理定律在昨天、今天和明天是完全相同的。做同一个实验，无论什么时候做，结果都一样。

  \textbf{第 2 步：Noether 定理给出守恒量。}
  对于时间平移对称性，Noether 定理给出的守恒量是\textbf{哈密顿量} $H$：
  \[
  H = \sum_i p_i \dot{q}_i - L
  \]
  其中 $p_i = \partial L/\partial \dot{q}_i$ 是广义动量，$\dot{q}_i$ 是广义速度，$L$ 是拉格朗日量。

  \textbf{第 3 步：验证哈密顿量确实守恒。}
  对 $H$ 求时间导数，利用 Euler-Lagrange 方程（来自 R2）可以证明 $dH/dt = -\partial L/\partial t$。由于我们已经假设 $\partial L/\partial t = 0$（时间平移不变），所以：
  \[
  \frac{dH}{dt} = 0
  \]

  \textbf{第 4 步：将 $H$ 识别为总能量。}
  在标准力学系统中（$L = T - V$，动能减势能），$H = T + V$——正是总机械能。对于更一般的系统，我们直接将 $H$ 定义为系统的总能量 $E$。因此时间平移不变性 $\to$ 能量守恒。

  \textbf{具体例子}：一个无摩擦的摆。拉格朗日量 $L = \frac12 m\dot{\theta}^2 - mg\ell(1-\cos\theta)$ 不显含时间。因此摆的总能量 $E = \frac12 m\dot{\theta}^2 + mg\ell(1-\cos\theta)$ 守恒——动能和势能不断互相转换，但总和始终不变。
\end{derivationbox}

\section{动量守恒}

\begin{lawbox}[Momentum Conservation]
\[
\frac{d\mathbf{P}}{dt} = 0 \quad\text{（孤立系统的总动量不变）}
\]
局域形式：$\partial_\mu T^{\mu i} = 0$。

\textbf{前提}：Noether 定理 + 空间平移对称（物理定律在空间各处相同）

\textbf{变量}：
\begin{itemize}
  \item $\mathbf{P} = \sum_i m_i \mathbf{v}_i$ — 总动量（kg·m/s）。系统中所有粒子的质量乘速度的矢量和。动量是一个矢量——它有大小和方向
  \item $d\mathbf{P}/dt = 0$ — 总动量不随时间变化。系统内部可以有复杂的动量交换，但总动量始终保持不变
\end{itemize}
\end{lawbox}

\begin{derivationbox}[推导：空间平移不变 → 动量守恒]
  \textbf{第 1 步：什么是"空间平移不变"？}
  如果你把整个实验设备从北京搬到纽约，实验结果不会改变。这意味着物理定律不依赖于你身在何处——空间是均匀的。在拉格朗日表述中，这意味着 $L$ 在空间平移 $\mathbf{r} \to \mathbf{r} + \boldsymbol{\epsilon}$ 下不变。

  \textbf{第 2 步：Noether 定理给出守恒量。}
  对于空间平移对称性，Noether 定理给出的守恒量是总动量：
  \[
  \mathbf{P} = \sum_i \frac{\partial L}{\partial \dot{\mathbf{r}}_i} = \sum_i m_i \dot{\mathbf{r}}_i = \sum_i m_i \mathbf{v}_i
  \]

  \textbf{第 3 步：动量守恒 → 牛顿第三定律。}
  动量守恒是牛顿第三定律的根本原因。对于两个相互作用的粒子：
  \[
  \frac{d}{dt}(\mathbf{p}_1 + \mathbf{p}_2) = 0 \;\implies\; \frac{d\mathbf{p}_1}{dt} + \frac{d\mathbf{p}_2}{dt} = 0
  \]
  由牛顿第二定律 $\mathbf{F} = d\mathbf{p}/dt$：
  \[
  \mathbf{F}_{12} + \mathbf{F}_{21} = 0 \;\implies\; \mathbf{F}_{12} = -\mathbf{F}_{21}
  \]
  这就是牛顿第三定律——作用力等于反作用力。它不过是动量守恒在两体系统中的表现形式。

  \textbf{具体例子}：一个静止的溜冰者推另一个静止的溜冰者。两人都向后滑，但总动量始终为零（$0 = m_1v_1 + m_2v_2$，所以 $v_2 = -(m_1/m_2)v_1$——质量大的滑得慢，质量小的滑得快）。
\end{derivationbox}

\section{角动量守恒}

\begin{lawbox}[Angular Momentum Conservation]
\[
\frac{d\mathbf{L}}{dt} = 0 \quad\text{（在有心力场中，总角动量不变）}
\]
\textbf{前提}：Noether 定理 + 旋转对称（物理定律在所有方向上相同）

\textbf{变量}：
\begin{itemize}
  \item $\mathbf{L} = \sum_i \mathbf{r}_i \times \mathbf{p}_i$ — 总角动量（kg·m$^2$/s）。角动量衡量系统的"旋转量"。$\times$ 是叉积（cross product）：$\mathbf{L}$ 的方向由右手定则确定，大小 $|\mathbf{L}| = rp\sin\theta$ 其中 $\theta$ 是 $\mathbf{r}$ 与 $\mathbf{p}$ 的夹角
  \item $\mathbf{r}_i$ — 第 $i$ 个粒子相对参考点的位置矢量
  \item $\mathbf{p}_i = m_i\mathbf{v}_i$ — 第 $i$ 个粒子的动量
\end{itemize}
\end{lawbox}

\begin{derivationbox}[推导：旋转不变 → 角动量守恒]
  \textbf{第 1 步：什么是"旋转不变"？}
  你朝北做一个实验，然后转向朝东再做一次——结果相同。空间在所有方向上是等价的（各向同性）。在拉格朗日表述中，这意味着 $L$ 在空间旋转下不变。

  \textbf{第 2 步：Noether 定理给出守恒量。}
  对于旋转对称性，Noether 定理给出的守恒量是总角动量 $\mathbf{L} = \sum_i \mathbf{r}_i \times \mathbf{p}_i$。

  \textbf{第 3 步：有心力场的特殊情况。}
  在有心力场中（力的方向总是指向或背离一个固定中心，如太阳引力），力矩为零：
  \[
  \boldsymbol{\tau} = \frac{d\mathbf{L}}{dt} = \mathbf{r}\times\mathbf{F}
  \]
  如果 $\mathbf{F} \parallel \mathbf{r}$（力与位置矢量平行），则 $\mathbf{r}\times\mathbf{F} = 0$，因此 $d\mathbf{L}/dt = 0$。

  \textbf{具体例子}：花样滑冰旋转。滑冰者收拢手臂时，质量到旋转轴的距离 $r$ 减小。由于角动量 $L = mrv$ 守恒，$r$ 减小必然意味着 $v$ 增大——滑冰者旋转得更快。这就是角动量守恒的日常体现。同样的原理解释了为什么行星在近日点比远日点运动得快（Kepler 第二定律——相等时间扫过相等面积）。
\end{derivationbox}

\section{电荷守恒}

\begin{lawbox}[Charge Conservation]
\[
\frac{dQ}{dt} = 0 \quad\text{（孤立系统的总电荷不变）}
\]
局域形式（连续性方程）：$\partial_\mu J^\mu = 0$，展开为 $\partial\rho/\partial t + \nabla\cdot\mathbf{J} = 0$。

\textbf{前提}：Noether 定理 + $U(1)$ 全局规范对称（R3）

\textbf{变量}：
\begin{itemize}
  \item $Q$ — 总电荷（C，库仑）。孤立系统的总电荷永远不变
  \item $\rho$ — 电荷密度（C/m$^3$）。每立方米中的电荷量
  \item $\mathbf{J}$ — 电流密度（A/m$^2$）。每平方米截面积流过的电流
  \item $\partial\rho/\partial t + \nabla\cdot\mathbf{J} = 0$ — 连续性方程。"某点电荷密度的减少 = 电荷从该点流出的速率"。这不是一个额外的物理定律——它是 $U(1)$ 规范对称性的必然结果
\end{itemize}
\end{lawbox}

\begin{derivationbox}[{推导：$U(1)$ 规范对称 → 电荷守恒}]
  \textbf{第 1 步：什么是 $U(1)$ 规范对称？}
  在量子力学中，物理可观测量只依赖于波函数的模方 $|\psi|^2$。因此，将波函数乘以一个复数相位因子 $e^{i\alpha}$（$\alpha$ 是常数）不会改变任何可观测物理量：$\psi \to e^{i\alpha}\psi$，$|\psi|^2$ 不变。这就是 $U(1)$ 全局对称性——"全局"意味着 $\alpha$ 在时空所有点都相同。

  \textbf{第 2 步：Noether 定理给出守恒流。}
  对于 Dirac 费米子场 $\psi$，其拉格朗日密度 $\mathcal{L} = \bar{\psi}(i\gamma^\mu\partial_\mu - m)\psi$ 在全局 $U(1)$ 变换下不变。Noether 定理给出守恒流：
  \[
  J^\mu = e\,\bar{\psi}\gamma^\mu\psi
  \]
  其中 $e$ 是基本电荷（$1.602 \times 10^{-19}$ C），$\bar{\psi} = \psi^\dagger\gamma^0$，$\gamma^\mu$ 是 Dirac 矩阵。守恒条件为 $\partial_\mu J^\mu = 0$。

  \textbf{第 3 步：从守恒流到电荷守恒。}
  对守恒流进行空间积分。$J^0 = \rho$（电荷密度）。由 $\partial_\mu J^\mu = 0$：
  \[
  \frac{\partial\rho}{\partial t} + \nabla\cdot\mathbf{J} = 0
  \]
  在一个体积 $V$ 上积分，用 Gauss 定理将散度项转换为面积分：
  \[
  \frac{d}{dt}\int_V \rho\,d^3x = -\oint_{\partial V}\mathbf{J}\cdot d\mathbf{A}
  \]
  即"体积内电荷的变化率 = -流出体积的净电流"。对于孤立系统（表面无电流进出），$dQ/dt = 0$。

  \textbf{具体例子}：你可以创造电子-正电子对（$e^- + e^+$），它们各自的电荷是 $-e$ 和 $+e$，总电荷为 $0$。电子和正电子的电荷精确相消——总电荷始终为零。电荷永远守恒，因为电磁相互作用的 $U(1)$ 规范结构保证了它。
\end{derivationbox}

\section{连续性方程——所有守恒律的局域形式}

上文中反复出现的 $\partial_\mu J^\mu = 0$ 形式是所有守恒律的统一语言。它的物理意义是：\textbf{守恒不是远距离的"神秘同步"，而是每一点的"流入必须等于流出加上积累"}。如果某个体积内的电荷减少了，一定是因为电荷流出了那个体积——不可能凭空消失。

\section{推导关系总图}

\begin{center}
\begin{tikzpicture}[
  node distance=1.0cm,
  kernel/.style={rectangle,draw,fill=principleblue!15,rounded corners,text width=2.2cm,align=center,font=\footnotesize\bfseries},
  sym/.style={rectangle,draw,fill=principleblue!5,rounded corners,text width=2.5cm,align=center,font=\footnotesize},
  cons/.style={rectangle,draw,fill=deriveorange!15,rounded corners,text width=2.3cm,align=center,font=\footnotesize},
  arrow/.style={->,>=Stealth,thick,deriveorange}
]
  \node[kernel] (r2) {R2\\$\delta S = 0$};
  \node[sym,below left=0.8cm and -0.5cm of r2] (tt) {时间平移对称\\$L$ 不依赖于 $t$};
  \node[sym,below=0.8cm of r2] (st) {空间平移对称\\$L$ 不依赖于 $\mathbf{r}$};
  \node[sym,below right=0.8cm and -0.5cm of r2] (ro) {旋转对称\\$L$ 在旋转下不变};
  \node[cons,below=of tt] (ene) {能量 $E$ 守恒};
  \node[cons,below=of st] (mom) {动量 $\mathbf{p}$ 守恒};
  \node[cons,below=of ro] (ang) {角动量 $\mathbf{L}$ 守恒};
  \node[sym,right=1.5cm of ro] (u1) {$U(1)$ 规范对称\\$\psi\to e^{i\alpha}\psi$};
  \node[cons,below=of u1] (chg) {电荷 $Q$ 守恒};
  \draw[arrow] (r2) -- (tt); \draw[arrow] (r2) -- (st);
  \draw[arrow] (r2) -- (ro); \draw[arrow] (tt) -- (ene);
  \draw[arrow] (st) -- (mom); \draw[arrow] (ro) -- (ang);
  \draw[arrow] (r2) -- (u1); \draw[arrow] (u1) -- (chg);
\end{tikzpicture}
\end{center}

Noether 定理是物理学中最深刻的联系之一：守恒律不是偶然的经验事实，而是宇宙几何对称性的直接表现。如果明天物理定律突然改变（时间平移对称破缺），能量就不再守恒。今天我们能放心地使用能量守恒来设计发动机、分析电路、理解恒星——是因为我们相信，或者说实验反复验证了，宇宙的底层对称性。

参考文献：[NOETHER] (1918), [GOLD] Ch.\ 12, [JACK] §5.2
